\documentclass[11pt]{article}
\usepackage{fontspec}
\usepackage[french]{babel}
\usepackage{amsmath}
\usepackage{amssymb}
\usepackage{amsthm}
\usepackage{mathrsfs}
\usepackage{tikz-cd}
\usepackage{enumitem}

\newcommand{\printedpage}[1]{\marginpar{\footnotesize p.~#1}}
\newcommand{\msp}[1]{\marginpar{\footnotesize [#1]}}
% printed notation: the category C' enclosed in a V-shaped wedge with the arrow name under it
\newcommand{\Vcat}[2]{\underset{#2}{\underset{\vee}{#1}}}

\begin{document}

% ===================== PDF page 1 =====================
\printedpage{--}
% no page number printed (title page)
% head: (none)

\begin{center}
\vspace*{4cm}
{\Large\textsf{LES DÉRIVATEURS}}

\bigskip\bigskip

\textbf{ALEXANDRE GROTHENDIECK}

\bigskip\bigskip\bigskip

\textbf{Chapitre II}

\textbf{Cofinalité (à droite et à gauche)}

\textbf{(préliminaire à la cofinalité cohomologique)}
\end{center}

\clearpage

% ===================== PDF page 2 =====================
\printedpage{--}
% no page number printed (editorial note)
% head: (none)

\vspace*{6cm}

\noindent Ce texte a été déchiffré et transcrit en \LaTeXe{} par M. Künzer. Il a été édité par M. Künzer, J. Malgoire, et G. Maltsiniotis. La transcription est aussi fidèle que possible au manuscrit. Pour les quelques corrections évidentes, ou rares commentaires des éditeurs, ainsi que pour la numérotation originale des pages du manuscrit, les caractères de machine à écrire \texttt{[typewriter]} entre crochets sont utilisés. Un point d'interrogation entre crochets signifie que l'on n'est pas sûr du mot qui précède.

\noindent Cette édition est provisoire. Les remarques, commentaires et corrections sont bienvenus. Envoyer un message à :

\noindent\textbf{maltsin@math.jussieu.fr}

\hfill G. Maltsiniotis\hspace*{2cm}

\clearpage

% ===================== PDF page 3 =====================
\printedpage{1}
% head: A. Grothendieck : Dérivateurs, Ch. II, Édition des Universités de Montpellier II et Paris VII
% foot: Version 8 juillet 2001

\msp{page 1}

\section*{1\quad Cofinalité cohomologique}

Soit
\[ \alpha : C' \longrightarrow C \]
un foncteur de (petites) catégories, $\mathcal{M}$ une catégorie ayant les $\varprojlim$ pertinentes, $F$ un préfaisceau (donc contravariant en $i$) sur $C$ à valeurs dans $\mathcal{M}$. On a un homomorphisme canonique dans $\mathcal{M}$
\[ \varprojlim_{C} F \longrightarrow \varprojlim_{C'} \underbrace{\alpha^*(F)}_{F\circ\alpha} \tag*{(1)} \]
et on veut des conditions sur $\alpha$ pour que ce soit toujours un isomorphisme (\footnote{On dit alors que $\alpha$ est \textit{cofinal,} cf. SGA 4, I 8.1, ou 3.3, p. 45.}). On est ramené aussitôt au cas où $\mathcal{M} = (\mathrm{Ens})$. Disons qu'une catégorie $I$ est \textit{faiblement filtrante à gauche} si $I \neq \emptyset$ et si $\forall\, i', i'' \in \operatorname{Ob} I$, $\exists\, i \in \operatorname{Ob} I$ et des flèches
\[
\begin{tikzcd}[row sep=small, column sep=small]
 & i' \\
i \arrow[ur] \arrow[dr] & \\
 & i''.
\end{tikzcd}
\]
Cela signifie aussi que $\operatorname{Ob} I$ en tant qu'ensemble ordonné est filtrant à gauche.

\textbf{Proposition.} \textit{Pour que} (1) \textit{soit toujours un isomorphisme} ($\alpha$ \textit{`anodin pour} $\varprojlim$ \textit{'), il suffit que pour tout} $x_0 \in \operatorname{Ob} C$\textit{, la catégorie} $x_0\backslash C'$ \textit{des couples} $(x', u)$\textit{, avec} $x' \in \operatorname{Ob} C'$\textit{,} $u : \alpha(x') \longrightarrow x$\textit{, soit faiblement filtrante à gauche} (\footnote{Cf. amélioration des remarques p. 5.}, \footnote{Cf. 3.2, p. 43-44, pour conditions nécessaires et suffisantes de cofinalité.}). % [sic] u : \alpha(x') \longrightarrow x -> x_0 in place of x, the category is x_0\backslash C'

Ceci n'implique pas la conclusion plus forte
\[ F \underbrace{\xrightarrow{\;\sim\;}}_{\text{adj.}} \alpha_*\alpha^* F \]
que le morphisme d'adjonction canonique soit un isomorphisme. Pour qu'il soit ainsi, il suffit que $\forall\, x_0 \in C$, le foncteur $\alpha/x_0 : C'/x_0 \longrightarrow C/x_0$ induit par $\alpha$ soit `anodin pour les $\varprojlim$ ', ce qui signifie aussi que pour toute flèche $u : x_1 \longrightarrow x_0$ dans $C$, la catégorie $\Vcat{C'}{u}$ formée des triples $(x', \xi, \eta)$ avec $x' \in \operatorname{Ob} C'$ et $\xi : x_1 \longrightarrow \alpha(x')$, $\eta : \alpha(x') \longrightarrow x_0$ rendant commutatif
% printed: C' inside a V-shaped wedge, with u printed below the wedge
\msp{page 2}
le diagramme
\printedpage{2}
% head: A. Grothendieck : Dérivateurs, Ch. II, Édition des Universités de Montpellier II et Paris VII
% foot: Version 8 juillet 2001
\[
\begin{tikzcd}
x_1 \arrow[rr, "u"] \arrow[dr, "\xi"'] & & x_0 \\
 & \alpha(x'), \arrow[ur, "\eta"'] &
\end{tikzcd}
\]
est faiblement filtrante à gauche.

Plus généralement, considérons
\[
\begin{tikzcd}
X' = C' \arrow[rr, "\alpha"] \arrow[dr, "f'"'] & & C = X \arrow[dl, "f"] \\
 & S &
\end{tikzcd}
\]
et posons nous la question quand l'homomorphisme fonctoriel, par rapport à un préfaisceau $F$ variable sur $X$,
\[ f_*(F) \longrightarrow f'_*(\alpha^*(F)) \]
est toujours un isomorphisme. (Le cas précédent est celui où $S = X$, $f = \mathrm{id}_X$.) On trouve qu'il suffit que $\forall\, s_0 \in \operatorname{Ob} S$, le foncteur
\[ \alpha/s_0 : X'/s_0 \longrightarrow X/s_0 \]
soit `anodin pour les $\varprojlim$ ', et pour ceci, il suffit que pour \texttt{[tout]} objet $(x_0, u_0)$ de $X/s_0$ (où
\[ f(x_0) \xrightarrow{u_0} s_0) \]
la catégorie $(x_0, u_0)\backslash(X'/s_0)$ soit faiblement filtrante à gauche.

\msp{page 3}

Cette catégorie s'interprète comme celle des triples $(x', \xi, u)$ avec
\[
\left\{
\begin{aligned}
& x' \in \operatorname{Ob} C' \\
& u : g(x') = f(\alpha(x')) \longrightarrow s_0 \\
& \xi : x_0 \longrightarrow \alpha(x')
\end{aligned}
\right.
\]
\texttt{[plutôt} $f'(x')$ \texttt{au lieu de} $g(x')$\texttt{]} de façon à rendre commutatif le diagramme
\[
\begin{tikzcd}
f(x_0) \arrow[rr, "u_0"] \arrow[dr, "f(\xi)"'] & & s_0 \\
 & f(\alpha(x')). \arrow[ur, "u"'] &
\end{tikzcd}
\]
Je voudrais dégager des variantes de ces énoncés dans le cadre cohomologique, les foncteurs $\varprojlim$ étant remplacés par des $R\Gamma$ (de préfaisceaux abéliens p.ex. - mais on pense aussi aux `cas non commutatifs', définis par des dérivateurs notamment.)

\clearpage

% ===================== PDF page 5 =====================
\printedpage{3}
% head: A. Grothendieck : Dérivateurs, Ch. II, Édition des Universités de Montpellier II et Paris VII
% foot: Version 8 juillet 2001

Si $I$ est une catégorie \textit{ordonnée finie}, la condition \texttt{[faiblement]} filtrante à gauche équivaut à l'existence d'un objet initial, ce qui implique que $I$ est contractile. Dans les énoncés précedents, il se pourrait que la bonne philosophie soit de remplacer partout `faiblement filtrante à gauche' par `contractile'. Si p.ex. $C$ est la catégorie finale, alors remplacer la condition que $C'$ soit filtrante à gauche, par l'acyclicité. C'est en tous cas en accord avec l'esprit de l'axiome $\mathrm{Der}_5$ (cas où $J$ est la catégorie finale). En effet, dans le cas où $C = e$, la condition que $\alpha$ soit anodin pour les $\varprojlim$ équivaut visiblement à dire que le foncteur $\alpha^* : \mathcal{M} \longrightarrow \underline{\mathrm{Hom}}(C'^{o}, \mathcal{M})$, associant à tout $F \in \mathcal{M}$ le préfaisceau constant de valeur $F$ % [sic] précedents -> précédents, missing accent on the second e
\msp{page 4}
sur $C'$, donne un morphisme d'adjonction
\[ \varprojlim_{e} F = F \longrightarrow \alpha_*(\alpha^*(F)) = \varprojlim_{C'} \alpha^* F \]
qui est un isomorphisme. Et l'axiome $\mathrm{Der}_5$ assure qu'il doit en être ainsi si $C'$ est contractile.

La condition filtrante dans la proposition est d'ailleurs une condition \textit{suffisante} d'anodinité, mais nullement nécessaire. Si p.ex. $C$ est la catégorie finale, alors la condition d'anodinité pour $C' \longrightarrow C$ est simplement que $C'$ soit connexe non vide, i.e. $\pi_0(C') = \{0\}$. (Condition qui visiblement néglige les $\pi_i$ supérieurs.)

Soit \fbox{$\alpha : X' \longrightarrow X$} un foncteur de petites catégories. La condition suggérée par la proposition, que les $x_0\backslash X'$ soient contractiles, signifie que $\alpha$ est une hot-équivalence co-stricte, i.e. que $\alpha^{o} : X'^{o} \longrightarrow X^{o}$ est une hot-équivalence `localement sur $X^{o}$', donc aussi globalement. Si $\alpha$
\msp{page 5}
lui-même était une hot-équivalence stricte, au lieu de $\alpha^{o}$, alors on aurait par $\mathrm{Der}_5$, que pour $F \in \mathbf{D}(X)$, $F \xrightarrow{\;\sim\;} \alpha_*\alpha^* F$ ($\mathbf{D}$ dérivateur à gauche), ce qui est bien plus fort que $\Gamma_{\mathbf{D}}(X, F) \xrightarrow[\sim]{\alpha^*} \Gamma_{\mathbf{D}}(X', \alpha^* F)$. La condition duale implique (si $\mathbf{D}$ \texttt{[est un]} dérivateur à droite) $\alpha_!\alpha^* F \xrightarrow{\;\sim\;} F$. Mais est-il raisonnable de s'attendre qu'elle implique aussi la condition `du côté cohomologie, au lieu de l'homologie',
\[ \Gamma_{\mathbf{D}}(X, F) \xrightarrow{\;\sim\;} \Gamma_{\mathbf{D}}(X', \alpha^*(F)) \ \ ? \]

Le premier cas, qui devrait emporter conviction, est celui du dérivateur DerAb, i.e. des catégories dérivées des faisceaux abéliens :

\textbf{Question.} \textit{Soit} $\alpha : X' \longrightarrow X$ \textit{une} hot-\textit{équivalence co-stricte,} $F$ \textit{un préfaisceau abélien sur} $X$\textit{. A-t-on alors} $H^i(X, F) \xrightarrow[\sim]{\alpha^*} H^i(X', \alpha^*(F))$ \textit{pour tout} $i \in \mathbf{Z}$ \textit{?}

Si $X$ admet un objet initial $\emptyset$, alors $\emptyset\backslash X' \simeq X'$, donc $X'$ contractile, donc c'est OK pour $F$ constant.

\textbf{Remarque.} Je m'aperçois que dans la proposition 1, au lieu de la condition que les $x\backslash C'$ soient faiblement filtrantes à gauche, il suffirait qu'elles soient non vides et connexes, i.e.

\end{document}
